\begin{ExerciseBlock}
\ExerciseHeading{07}{Rotierende Massen mit Feder}{two_particle_systems_rotating_masses_and_spring}
Zwei Massen mit $m_1 = 5\,$kg und $m_2 = 1\,$kg sind durch eine Feder mit der Federkonstanten $D=0,5\,$N/cm verbunden.
Im Ruhezustand haben die Massen einen Abstand von $r_0 = 1\,$m. \\

\begin{center}
	\CachedTikzFigure{content/resources/chapters/two_particle_systems/tikz/rotating_masses_and_spring.tex}
\end{center}

\begin{enumerate}
	\item Berechnen Sie den Abstand der beiden Massen zum gemeinsamen Schwerpunkt.

	\item Nun beginnt dieses System mit $\omega = 5 \, \text{s}^{-1}$ um seinen Schwerpunkt zu rotieren. Berechnen Sie den Abstand $r_{rot}$ der Massen, der sich in der Folge einstellt.

\end{enumerate}
\end{ExerciseBlock}
\UseExerciseSolution{two_particle_systems_rotating_masses_and_spring}
